% ============================================================
%  第二卷 第七章 光的传播（§53–§61 + 习题）
%  底本：高教社中文版；OCR 错误已修正，脚注文本待从纸版补录
%
%  OCR 错误修正清单：
%  1. §53 OCR 块1/块2 拼接断句（行6887/6889）："……那样简单的形式.但是
%     /在本质上……" 合并为一段一句；
%  2. §53 跨页断句（行6936/6938）："正如程函方/程一样" 接合为同段；
%  3. "引人"→"引入"，"所 $\ngeq$ 生"→"所产生"，"因 $\mathfrak{F}_1$ 对于"
%     →"因而对于"，"$\bar\kappa$ 消去"→"从而消去"，"$\bar{\kappa_{\bar n}}$
%     比起"→"因而比起"，"公式成 $\yen 0$"→"公式成立"；
%  4. §53 (53.9) $\delta=A\omega$→$\mathcal{E}=A\omega$；"用波矢k代替动量"
%     后的 $\pmb{w}$→$\omega$；$\pmb{\pi}$→$\pi$；
%  5. §54 "设$^{ac}$和bd"→"设 $ac$ 和 $bd$"；曲率中心/线元字母恢复斜体；
%  6. §55–§56 OCR 大量 \pmb/\mathbf/\Delta 乱码矢量记号统一为 \boldsymbol；
%  7. §56 习题2 $p^2=\ell^2/c^2-m^2c^2$→$p^2=\mathcal{E}^2/c^2-m^2c^2$
%     （\ell 为 \mathcal{E} 之误，已对照原书 p.158）；
%  8. §56 习题内编号公式 (1)(2)(3) 按原书用手工 \tag（节内自动编号规则不适用于题内编号）；
%  9. §59 "由引可见"→"由此可见"（原书排印之误）；§60 "距离为 $D_q$ 我们用"
%     补句号；§60 "比 $D_g$ 和 $D_p$ 都小"→$D_q$；
% 10. §58 习题 "最小值 $\sim\sqrt{\lambda}l$"→$\sqrt{\lambda l}$
%     （按物理量纲修正，原书根号仅罩住 $\lambda$）；
% 11. §60 替换式中 $\left(\frac{1}{D_p}+\frac{1}{D_q}\right)y$ 的 $y$→$z$
%     （与 (60.2)、对 z 积分一致，原书误排 y，已对照原书 p.171）；
% 12. §61 习题1 OCR 漏句"朝两边增加，强度历经一系列高度迅速减小的极大值.
%     相邻的极大值被处于"按原书 p.175 补回；
% 13. §61 (61.5) 后乱码"$\vec{x}\cdot\mp\cdot q\neq 0$"→"对于 $q\neq 0$"；
%     习题2 lim 下标 "$N\ \infty$"→$N\to\infty$；
% 14. (60.5) 中积分下限 $-\omega$→$-w$；各处句末缺句号补齐；
% 15. 粗体矢量：OCR 将 \boldsymbol 丢失/误作 \pmb、\mathbf、\Delta 者按原书恢复。
%
%  脚注圈码 ①②③ 按规范保留在正文中；原书相应脚注文本未转录。
%  §59 习题与 §61 习题2 的脚注在 OCR 正文中带有全文，已整段注释保留于原位，
%  供脚注代理转录时取用。
%
%  脚注已转录（2026-10-10）：10 处圈码全部替换为 \footnote{...}
%  （转录自原书扫描页 p-165/168/172/173/178/182/187/190，与 MinerU OCR
%  page_footnote 块逐字核对）。§61 习题2 脚注中"由此可见……（见 §28 的
%  脚注）"一句原书属脚注末句，已从正文移入脚注。§59 习题脚注（艾里函数）
%  与草稿注释一致。
%
%  遗留问题：
%  * §53 (53.1) 中 $\mathrm{Re}$ 说明文字与原书一致；
%  * §59 艾里函数脚注中"（见本教程第三卷 §b）"之 §b 照原书照排。
% ============================================================
\chapter{光的传播}

\section{几何光学}

平面波的特征是它的传播方向和振幅处处相同.任意的电磁波当然没有这种特性.

然而，在许多情况下，电磁波虽然不是平面波，但具有如下特性，即它们在空间的每一个小区域内可以当做是平面波.为了满足这个要求，显然，波的振幅和方向在与波长同数量级的距离内必须是几乎不变的.

如果这个条件被满足了，我们可以引入一个所谓波面，它是这样一个面，在该面上的所有点，波的相位（在给定时刻）都是一样的（平面波的波面显然是一个与光的传播方向相垂直的平面）.在空间的每一个小区域内，我们可以说光的传播方向垂直于波面.这样，我们可以引入光线这个概念，光线是其上每一点的切线都同光的传播方向相合的线.

用这种方法研究波的传播规律属于几何光学的范畴.因此，几何光学将波的传播，特别是光的传播，当做光线的传播，因而完全同它的波的特性脱离了关系.换句话说，几何光学相当于波长 $\lambda\to 0$ 的极限情形.

我们现在来求几何光学的基本方程——决定光线方向的方程.设 $f$ 是描写波场的任意一个量（$\boldsymbol{E}$ 或 $\boldsymbol{H}$ 的任意一个分量）.在单色平面波中，$f$ 有下面的形式：
\begin{equation}
  f=a\exp[\mathrm{i}(\boldsymbol{k}\cdot\boldsymbol{r}-\omega t+\alpha)]
  =a\exp[\mathrm{i}(-k_ix^i+\alpha)]
\end{equation}
（我们略去了符号 $\mathrm{Re}$；以后遇到的所有式子都理解为取其实数部分）.

我们将场的表达式写成
\begin{equation}
  f=a\,\mathrm{e}^{\mathrm{i}\psi}.
\end{equation}

对于波不是平面的，但几何光学可以适用的情况，振幅 $a$ 一般来说是坐标与时间的函数，而相位 $\psi$，也叫做程函，没有像（53.1）那样简单的形式.但是在本质上，$\psi$ 是一个很大的量.这可以直接从下面的事实看出来，即当我们移动一个波长时，$\psi$ 就变动 $2\pi$，而几何光学是与极限 $\lambda\to 0$ 相应的.

在一个小的空间区域中和一个小的时间间隔内，程函可以展开为级数；如果只精确到第一级，则我们有
\begin{equation*}
  \psi=\psi_0+\boldsymbol{r}\cdot\frac{\partial\psi}{\partial\boldsymbol{r}}
  +t\frac{\partial\psi}{\partial t}
\end{equation*}
（坐标及时间原点都选择在所研究的空间区域及时间间隔内；导数在原点取值）.将上式与（53.1）式相比较，得到
\begin{equation}
  \boldsymbol{k}=\frac{\partial\psi}{\partial\boldsymbol{r}}
  \equiv\operatorname{grad}\psi,\qquad
  \omega=-\frac{\partial\psi}{\partial t},
\end{equation}
这相应于下面的事实，即在空间的每一个小区域内（且在每个小的时间间隔内），波可以当做是平面的.用四维空间形式，（53.3）式可以写成
\begin{equation}
  k_i=-\frac{\partial\psi}{\partial x^i},
\end{equation}
其中 $k_i$ 是波四维矢量.

在 §48 中我们已经看出，四维矢量 $k^i$ 的分量之间有 $k_ik^i=0$ 的关系.将（53.4）代入，得到
\begin{equation}
  \frac{\partial\psi}{\partial x_i}\frac{\partial\psi}{\partial x^i}=0.
\end{equation}

这个方程称为程函方程，它是几何光学的基本方程.

在波动方程内作直接的极限过渡 $\lambda\to 0$，也可以导出程函方程.场 $f$ 满足波动方程
\begin{equation*}
  \frac{\partial^2 f}{\partial x_i\partial x^i}=0.
\end{equation*}
将 $f=a\mathrm{e}^{\mathrm{i}\psi}$ 代入，得到
\begin{equation}
  \frac{\partial^2 a}{\partial x_i\partial x^i}\mathrm{e}^{\mathrm{i}\psi}
  +2\mathrm{i}\frac{\partial a}{\partial x_i}\frac{\partial\psi}{\partial x^i}\mathrm{e}^{\mathrm{i}\psi}
  +\mathrm{i}f\frac{\partial^2\psi}{\partial x_i\partial x^i}
  -\frac{\partial\psi}{\partial x_i}\frac{\partial\psi}{\partial x^i}f=0.
\end{equation}

但是上面已经指出，程函 $\psi$ 是一个很大的量；因此前三项与第四项相比可以略去，于是我们重又得到方程（53.5）.

在这里，我们还将求出一系列关系.虽然将这些关系应用于光在真空中的传播时只能导出一些显而易见的结果，不过，它们还是重要的.因为就其一般形式而言，这些推导也适用于光在物质介质中的传播.

从程函方程的形式，得到一个在几何光学与实物粒子力学间的惊人的相似性.粒子运动决定于哈密顿-雅可比方程（16.11）.这个方程，正如程函方程一样，是一阶偏导数的、二次的方程.我们知道，作用量 $S$ 与粒子的动量 $\boldsymbol{p}$ 及哈密顿量 $\mathcal{H}$ 有下列关系：
\begin{equation*}
  \boldsymbol{p}=\frac{\partial S}{\partial\boldsymbol{r}},\qquad
  \mathcal{H}=-\frac{\partial S}{\partial t}.
\end{equation*}
将这些公式与（53.3）比较，我们看出，波矢在几何光学中所扮演的角色正如粒子的动量在力学中所扮演的角色一样，而频率所扮演的角色是哈密顿量，即粒子的能量.波矢的绝对值 $k$ 同频率的关系为 $k=\omega/c$.我们看出，这个关系正与一个质量为零而速度等于光速的粒子的动量与能量的关系 $p=\mathcal{E}/c$ 相似.

对于质点，哈密顿方程
\begin{equation*}
  \dot{\boldsymbol{p}}=-\frac{\partial\mathcal{H}}{\partial\boldsymbol{r}},\qquad
  \boldsymbol{v}=\dot{\boldsymbol{r}}=\frac{\partial\mathcal{H}}{\partial\boldsymbol{p}}
\end{equation*}
成立.鉴于上述的相似性，我们可以直接写出光线的类似方程
\begin{equation}
  \dot{\boldsymbol{k}}=-\frac{\partial\omega}{\partial\boldsymbol{r}},\qquad
  \dot{\boldsymbol{r}}=\frac{\partial\omega}{\partial\boldsymbol{k}}.
\end{equation}

在真空中，$\omega=ck$，因而 $\dot{\boldsymbol{k}}=\boldsymbol{0}$，$\boldsymbol{v}=c\boldsymbol{n}$（$\boldsymbol{n}$ 是沿着波的传播方向的单位矢量）；就是说，在真空中，光线是直线，光沿着这条直线以速度 $c$ 传播.

下面的考虑，更加清楚地说明了波的波矢与粒子的动量的相似性.我们来考虑一种波，它是许多频率不同的单色波在某一小间隔内的叠加，这个波只占据空间的某个有限区域（这就是所谓波包）.我们利用公式（32.6）及能量动量张量（48.15）来计算这个波的场的四维动量（对于每一个单色分量）.用某些平均值代替这个公式中的 $k^i$，我们得到如下形式的表达式
\begin{equation}
  P^i=Ak^i,
\end{equation}
这里四维矢量 $P^i$ 与 $k^i$ 的比例常数 $A$ 是一标量.写成三维形式，这个关系给出
\begin{equation}
  \boldsymbol{P}=A\boldsymbol{k},\qquad \mathcal{E}=A\omega.
\end{equation}

由此可见，当我们从一个参考系过渡到另一个参考系时，波包的动量和能量像波矢和频率一样变换.

再往前追溯这种相似，我们可以为几何光学建立一个原理，这个原理与力学中的最小作用量原理相似.然而，不能将它写成如 $\delta\int L\,\mathrm{d}t=0$ 那样的哈密顿形式，因为对于光线，不可能引入像粒子的拉格朗日量那样的函数.事实上，粒子的拉格朗日量 $L$ 与哈密顿量 $\mathcal{H}$ 有 $L=\boldsymbol{p}\cdot\partial\mathcal{H}/\partial\boldsymbol{p}-\mathcal{H}$ 的关系.如果用频率 $\omega$ 代替哈密顿量，用波矢 $\boldsymbol{k}$ 代替动量，那么，光学中的拉格朗日量就应该写成
$\boldsymbol{k}\cdot\partial\omega/\partial\boldsymbol{k}-\omega$. 但是这个式子等于零，因为 $\omega=ck$.前面已经提到过，光线的传播同一个质量为零的粒子的运动相似，因此，对于光线引入拉格朗日量之不可能也是显而易见的.

如果波有一定的恒定频率 $\omega$，那么，它的场与时间的关系就为一个因子 $\mathrm{e}^{-\mathrm{i}\omega t}$ 所决定.因此，这样的波的程函就是
\begin{equation}
  \psi=-\omega t+\psi_0(x,y,z),
\end{equation}
式中，$\psi_0$ 仅是坐标的函数.程函方程（53.5）现在取下面的形式：
\begin{equation}
  (\operatorname{grad}\psi_0)^2=\frac{\omega^2}{c^2}.
\end{equation}

波面就是定值程函的曲面，亦即曲面族 $\psi_0(x,y,z)=\mathrm{const}$.光线本身在每一点都垂直于相应的波面；光线的方向为梯度 $\nabla\psi_0$ 所决定.

众所周知，如果能量是常数，粒子的最小作用量原理也可以写成所谓莫培督原理的形式：
\begin{equation*}
  \delta S=\delta\int\boldsymbol{p}\cdot\mathrm{d}\boldsymbol{l}=0,
\end{equation*}
这里的积分是沿着两点间的粒子的轨道而取的.在这个式子内，我们假设动量是能量与坐标的函数.光线的类似原理是费马原理.在这种情形下，我们将费马原理写成类似的形式：
\begin{equation}
  \delta\psi=\delta\int\boldsymbol{k}\cdot\mathrm{d}\boldsymbol{l}=0.
\end{equation}

在真空中，$\boldsymbol{k}=\dfrac{\omega}{c}\boldsymbol{n}$，我们得到（$\mathrm{d}\boldsymbol{l}\cdot\boldsymbol{n}=\mathrm{d}l$）
\begin{equation}
  \delta\int\mathrm{d}l=0,
\end{equation}
这与光沿直线传播相应.

\section{强度}

在几何光学中，光波可以当做一束光线.然而光线本身仅仅决定光在每一点的传播方向，还余下来一个问题，即光的强度在空间分布的问题.

在所考虑的光线束的波面上，我们取出一个无限小的面元.从微分几何知道，在每一个曲面的每一点上，有两个（一般说是不同的）主曲率半径.设 $ac$ 和 $bd$（图7）为主曲率圆的线元，它们在波面的给定面元上.通过 $a$ 和 $c$ 的光线相交于相应的曲率中心 $O_1$，而通过 $b$ 和 $d$ 的光线则相交于另一曲率中心 $O_2$.

\par\nopagebreak\noindent\begin{minipage}{\linewidth}\centering\includegraphics[width=6cm]{fig7.pdf}\\[0.3em]{\small 图 7}\end{minipage}\par\nopagebreak

当从 $O_1$ 和 $O_2$ 出发的光束的开放角一定时，弧 $ac$ 和 $bd$ 之长显然与曲率半径 $R_1$ 和 $R_2$ 成正比（即与 $O_1O$ 及 $O_2O$ 成正比）.面元的面积与长度 $ac$ 和 $bd$ 之积成正比，即与 $R_1R_2$ 成正比.换句话说，如果考虑由一束光线所构成的波面元，那么，当沿着光线移动时，面元的面积将与 $R_1R_2$ 成正比地改变.

另一方面，光的强度，即能流密度，与一定的光能量通过的曲面面积成反比.因此，我们得到一个结果，即光的强度是
\begin{equation}
  I=\frac{\mathrm{const}}{R_1R_2}.
\end{equation}

这个公式必须理解如下.在每一条光线上（图7中的 $AB$），存在着确定的点 $O_1$ 及 $O_2$，这两点是所有与光线相交的波面的曲率中心.$OO_1$ 与 $OO_2$ 是从 $O$ 点（波面与光线相交之点）到 $O_1$ 和 $O_2$ 的距离，即波面在 $O$ 点的曲率半径 $R_1$ 及 $R_2$.因此，公式（54.1）决定光沿着一条光线上的强度变化，表示为到该线上二定点之距离的函数.我们强调指出，这个公式不能用来比较同一波面上的不同点的强度.

因为强度为场的模的平方所决定，我们就可以写出场本身沿光线的变化如下：
\begin{equation}
  f=\frac{\mathrm{const}}{\sqrt{R_1R_2}}\,\mathrm{e}^{\mathrm{i}kR},
\end{equation}
式中，在相因子 $\mathrm{e}^{\mathrm{i}kR}$ 内，我们既可以用 $R_1$ 也可以用 $R_2$ 代替 $R$.对于一给定的光线，$\mathrm{e}^{\mathrm{i}kR_1}$ 和 $\mathrm{e}^{\mathrm{i}kR_2}$ 两个量只相差一个常量因子，因为差 $R_1-R_2$ 即两个曲率中心的距离，是一个常数.

如果波面的两个曲率半径相等，那么，（54.1）和（54.2）就有如下的形式：
\begin{equation}
  I=\frac{\mathrm{const}}{R^2},\qquad
  f=\frac{\mathrm{const}}{R}\,\mathrm{e}^{\mathrm{i}kR}.
\end{equation}

当光是从一个点源射出时，就对应于这种情形（波面都是同心球，$R$ 是从光源到波面的距离）.

从（54.1）式，我们看出，在 $R_1=0$，$R_2=0$ 各点，即在波面的曲率中心，强度变为无穷大.将这个结论应用到光线束，我们得到在一给定光线束内，光的强度，一般来说，在两个曲面上变为无穷大，这两个曲面就是波面的所有曲率中心构成的几何图形.这两个曲面称为焦散面.在光束的波面是球面的特殊情形下，两个焦散面就合为一点（焦点）.

必须指出，按照微分几何中熟知的曲面族的曲率中心几何构形的性质，光线与焦散面相切.

我们应当记着，对于凸波面，波面的曲率中心可以不在光线本身，而在其延长线上，在发出光线的光学系统之外.在这种情形下，我们说它是虚焦散面（或虚焦点）.这时，光的强度无论在哪里都不会变为无限大.

至于光的强度变为无限大，实际上应当理解为，强度在焦散面上变大了，但仍是有限的（见 §59 中的习题）.形式上，光强度为无限大就说明几何光学近似不能应用到焦散面附近.与此有关的还有下述事实，即相位沿着光线的变化可以由公式（54.2）来决定，但是只限于不包含与焦散面相切之点的一段.以后在 §59 中我们还要证明，实际上当光线穿过焦散面时，场的相位减少 $\pi/2$.这就是说，如果在光线上与焦散面第一次相交之前的那一段上，场与因子 $\mathrm{e}^{\mathrm{i}kx}$（$x$ 是沿光线的坐标）成正比，那么，穿过焦散面以后，场将与 $\mathrm{e}^{\mathrm{i}(kx-\pi/2)}$ 成正比.在第二个焦散面的切点的附近，又会发生同样的情形，在这点以后，场将与 $\mathrm{e}^{\mathrm{i}(kx-\pi)}$ 成正比\footnote{尽管公式（54.2）本身在焦散面附近不成立，场的相位改变在形式上相应于这个公式中 $R_1$ 或 $R_2$ 正负号的改变（即乘以 $\mathrm{e}^{\mathrm{i}\pi}$）.}.

\section{角程函}

如果一条光线在真空中行进，投射在一个透明物体上，当它从这个物体射出时，它的方向与原来的方向一般来说是不同的.这种方向的改变当然与物体的特性和物体的形状有关.但是，我们可以求得一些一般的规律，这些规律给出光线经过一个任意物体时方向的改变.在此，我们只假定几何光学可以应用到在我们所考虑的物体内行进的光线.依照惯例，光线所穿过的透明物体称为光学系统.

按照 §53 所指出的光线传播与粒子运动的相似性，这些普遍定律，对这样的粒子运动方向的改变也是有效的，这个粒子在真空中最初沿直线运动，以后经过某一电磁场，再从这个场中穿出来到真空中.但是，为了确定起见，我们以后总是说光线的传播.

我们在上节已经看出，（对于一定频率的光）描述光线传播的程函方程可以写成（53.11）的形式.从现在起，为便利起见，我们用 $\psi$ 代表被常数 $\omega/c$ 除过了的程函 $\psi_0$.这时，几何光学的基本方程有下列形式：
\begin{equation}
  (\nabla\psi)^2=1.
\end{equation}

该方程的每一个解描述一个确定的光线束，其中经过空间给定点的光线的方向由 $\psi$ 在该点的梯度决定.然而为了我们的目的，这种描述是不够的，因为我们正在寻找的，是决定任意光线（而不是一个确定的光线束）经过一个光学系统的路径的普遍关系.因此，我们必须应用程函的这样一个形式，它可以描述所有一般可能的光线，也就是经过空间任意一对点的光线.通常形式的程函 $\psi(\boldsymbol{r})$ 是经过点 $\boldsymbol{r}$ 的某束光线的相位.现在我们应当引入为两点坐标函数 $\psi(\boldsymbol{r},\boldsymbol{r}')$ 的程函（$\boldsymbol{r}$ 和 $\boldsymbol{r}'$ 是光线起点和终点的径矢）.对于每一对点 $\boldsymbol{r},\boldsymbol{r}'$ 都可以有光线经过这两点，而 $\psi(\boldsymbol{r},\boldsymbol{r}')$ 是光线上两点 $\boldsymbol{r}$ 和 $\boldsymbol{r}'$ 的相位差（或如一般所称的，光程长）.从现在起，我们把 $\boldsymbol{r}$ 和 $\boldsymbol{r}'$ 理解为透过光学系统之前和之后光线上两点的径矢.

如果在 $\psi(\boldsymbol{r},\boldsymbol{r}')$ 中，认为径矢之一，例如 $\boldsymbol{r}'$，是给定的，那么，$\psi$ 是 $\boldsymbol{r}$ 的函数，描写一个确定的光线束，即经过 $\boldsymbol{r}'$ 点的光线束.$\psi$ 应当满足方程（55.1），方程中的微分是对 $\boldsymbol{r}$ 的分量进行的.同理，如果 $\boldsymbol{r}$ 固定，我们又得到一个 $\psi(\boldsymbol{r},\boldsymbol{r}')$ 的方程，因此
\begin{equation}
  (\nabla\psi)^2=1,\qquad (\nabla'\psi)^2=1.
\end{equation}

从上节我们知道，光线的方向由它的相位的梯度所决定.既然 $\psi(\boldsymbol{r},\boldsymbol{r}')$ 是在点 $\boldsymbol{r}$ 及点 $\boldsymbol{r}'$ 的相位差，在 $\boldsymbol{r}'$ 点的光线的方向就由矢量 $\boldsymbol{n}'=\partial\psi/\partial\boldsymbol{r}'$ 所决定，而在 $\boldsymbol{r}$ 点的光线的方向，则由矢量 $\boldsymbol{n}=-\partial\psi/\partial\boldsymbol{r}$ 决定.从（55.2）式可以看出，$\boldsymbol{n}$ 及 $\boldsymbol{n}'$ 是单位矢量：
\begin{equation}
  n^2=n'^2=1.
\end{equation}

四个矢量 $\boldsymbol{r},\boldsymbol{r}',\boldsymbol{n},\boldsymbol{n}'$ 相互间有一定的关系，因为其中两个（$\boldsymbol{n}$ 和 $\boldsymbol{n}'$）是某一函数 $\psi$ 对另外两个（$\boldsymbol{r}$ 和 $\boldsymbol{r}'$）的导数.至于函数 $\psi$ 本身，它是满足附加条件方程（55.2）的.

为了求出 $\boldsymbol{n},\boldsymbol{n}',\boldsymbol{r},\boldsymbol{r}'$ 的关系，最便利的方法是用另外一个量代替 $\psi$，对这个量不加任何附带条件（即不要求它满足任何微分方程）.我们可以按照下面的方法去做.在函数 $\psi$ 内，独立变量是 $\boldsymbol{r}$ 和 $\boldsymbol{r}'$，因此，对于微分 $\mathrm{d}\psi$，我们有
\begin{equation*}
  \mathrm{d}\psi=\frac{\partial\psi}{\partial\boldsymbol{r}}\cdot\mathrm{d}\boldsymbol{r}
  +\frac{\partial\psi}{\partial\boldsymbol{r}'}\cdot\mathrm{d}\boldsymbol{r}'
  =-\boldsymbol{n}\cdot\mathrm{d}\boldsymbol{r}
  +\boldsymbol{n}'\cdot\mathrm{d}\boldsymbol{r}'.
\end{equation*}

现在我们作一个勒让德变换，从 $\boldsymbol{r},\boldsymbol{r}'$ 变换到新的独立变量 $\boldsymbol{n},\boldsymbol{n}'$；亦即我们写
\begin{equation*}
  \mathrm{d}\psi=-\mathrm{d}(\boldsymbol{n}\cdot\boldsymbol{r})
  +\boldsymbol{r}\cdot\mathrm{d}\boldsymbol{n}
  +\mathrm{d}(\boldsymbol{n}'\cdot\boldsymbol{r}')
  -\boldsymbol{r}'\cdot\mathrm{d}\boldsymbol{n}',
\end{equation*}
引入函数
\begin{equation}
  \chi=\boldsymbol{n}'\cdot\boldsymbol{r}'-\boldsymbol{n}\cdot\boldsymbol{r}-\psi,
\end{equation}
我们得到
\begin{equation}
  \mathrm{d}\chi=-\boldsymbol{r}\cdot\mathrm{d}\boldsymbol{n}
  +\boldsymbol{r}'\cdot\mathrm{d}\boldsymbol{n}'.
\end{equation}

函数 $\chi$ 称为角程函；从（55.5）式可以看出，角程函内的独立变量是 $\boldsymbol{n}$ 和 $\boldsymbol{n}'$，没有附带条件加在 $\chi$ 上.事实上，方程（55.3）现在仅仅表示关于独立变量的条件：在矢量 $\boldsymbol{n}$ 的三个分量 $n_x,n_y,n_z$ 中，只有两个是独立的（对于 $\boldsymbol{n}'$ 也有类似情况）.我们以后选取 $n_y,n_z,n_y',n_z'$ 为独立变量；于是，
\begin{equation*}
  n_x=\sqrt{1-n_y^2-n_z^2},\qquad
  n_x'=\sqrt{1-n_y'^2-n_z'^2}.
\end{equation*}

将这些式子代入
\begin{equation*}
  \mathrm{d}\chi=-x\,\mathrm{d}n_x-y\,\mathrm{d}n_y-z\,\mathrm{d}n_z
  +x'\,\mathrm{d}n_x'+y'\,\mathrm{d}n_y'+z'\,\mathrm{d}n_z',
\end{equation*}
我们得到微分 $\mathrm{d}\chi$ 的表达式
\begin{equation*}
\begin{aligned}
  \mathrm{d}\chi={}&-\left(y-\frac{n_y}{n_x}x\right)\mathrm{d}n_y
  -\left(z-\frac{n_z}{n_x}x\right)\mathrm{d}n_z\\
  &+\left(y'-\frac{n_y'}{n_x'}x'\right)\mathrm{d}n_y'
  +\left(z'-\frac{n_z'}{n_x'}x'\right)\mathrm{d}n_z'.
\end{aligned}
\end{equation*}

由此，我们最后便求得下面的方程：
\begin{equation}
\begin{aligned}
  y-\frac{n_y}{n_x}x&=-\frac{\partial\chi}{\partial n_y}, &
  z-\frac{n_z}{n_x}x&=-\frac{\partial\chi}{\partial n_z},\\
  y'-\frac{n_y'}{n_x'}x'&=\frac{\partial\chi}{\partial n_y'}, &
  z'-\frac{n_z'}{n_x'}x'&=\frac{\partial\chi}{\partial n_z'},
\end{aligned}
\end{equation}

这就是我们所求的 $\boldsymbol{n},\boldsymbol{n}',\boldsymbol{r},\boldsymbol{r}'$ 间的一般关系.函数 $\chi$ 描述光线所经过的物体的特性（或者对于带电粒子运动的情形，$\chi$ 描述场的特性）.

当 $\boldsymbol{n}$ 和 $\boldsymbol{n}'$ 之值固定时，（55.6）中每一对方程都表示直线.这些直线正是穿过光学系统前后的光线.因此，方程（55.6）直接决定光学系统两边光线的路径.

\section{窄束光线}

在研究光线束经过光学系统时，所有光线都相交于一点的光线束是值得特别注意的（这样的光线束称为同心光线束）.

经过光学系统后，同心光线束一般不再是同心的了，亦即经过物体后，光线不再会聚在一点了.仅仅在特殊情况下，光才会从一个发光点出发，穿过光学系统后，又会聚在一点（发光点的像）\footnote{交点可以在光线本身之上，也可以在它们的延长线上；取决于这两种不同的情形，像被称为实像或者虚像.}.

可以证明（见 §57），整个同心光线束经过光学系统后仍然保持为同心光线束的唯一情形是所谓全同成像，在此情形下，像与物体的差别仅在于平移、转动或镜像反射.

因此，没有光学系统能给予具有一定大小的物体以完全清晰的像，只有全同成像的情形是例外\footnote{这种成像可以借助平面镜实现.}.在全同成像以外的任何其他情形，一个有一定大小的物体只能产生近似清晰，而不是完全清晰的像.

一个同心光线束近似变换到另一个同心光线束的最重要的情况，是在一条（对该光学系统而言）特定的线附近行进的十分窄的（张角甚小）光线束.这条线称为光学系统的光轴.

然而必须注意，甚至无限窄的光线束（在三维空间内）在一般情况下也不是同心的；我们已经看到（图7），即使在这样的光束内，不同的光线也不会相交于同一点（这种现象称为像散）.只有两个主曲率半径相等的波面上的那些点是例外，这些点附近的波面上的小区域可以当做是球面，其相应的窄光线束是同心的.

我们来考虑一个具有轴对称的光学系统\footnote{可以证明，借助于在光轴近旁行进的窄光线束来处理的在一个非轴对称光学系统内成像的问题，可以化为在轴对称系统内成像，加上如此所得的像相对于物体的旋转.}.这个系统的对称轴同时也是它的光轴.沿着这个轴行进的光线束的波面也有轴对称性；正如我们所知道的，旋转面在它们与对称轴相交的那些点有相等的曲率半径.因此在这个方向行进的窄光线束保持为同心的.

为了得到定量的关系，以及借助窄光线束来决定经过一个轴对称光学系统的像的形成，我们利用通式（55.6）.首先当然要确定函数 $\chi$ 在所研究的情形中应取的形式.

因为光线束是窄的，而且在光轴的近旁行进，那么，每束光线的矢量 $\boldsymbol{n}$ 及 $\boldsymbol{n}'$ 差不多都沿着这个轴的方向.如果我们选择光轴作为 $x$ 轴，那么，分量 $n_y,n_z,n_y',n_z'$ 就比 1 小了.对于 $n_x,n_x'$ 而言，$n_x\approx 1$，而 $n_x'$ 则近似地等于 $+1$ 或 $-1$. 在第一种情形下，光线几乎沿着原来方向继续行进，穿过光学系统而进入其另一边的空间内.这种光学系统称为透镜.在第二种情形，光线改变方向到几乎相反的方向上；这种光学系统称为反射镜.

利用 $n_y,n_z,n_y',n_z'$ 值很小的性质，将角程函 $\chi(n_y,n_z,n_y',n_z')$ 展开为级数并只保留前几项.由于整个系统的轴对称性，$\chi$ 对于坐标系绕光轴的转动来说是不变量.由此可见，在 $\chi$ 的展开式中不可能有与矢量 $\boldsymbol{n}$ 和 $\boldsymbol{n}'$ 的 $y$ 分量和 $z$ 分量的一次幂成比例的一阶项；这样的项不可能具有所要求的不变性.具有所要求特性的二阶项是 $n^2,n'^2$ 和标积 $\boldsymbol{n}\cdot\boldsymbol{n}'$.因此，精确到二阶项为止的轴对称光学系统的角程函有下面的形式：
\begin{equation}
  \chi=\mathrm{const}+\frac{g}{2}(n_y^2+n_z^2)
  +f(n_yn_y'+n_zn_z')
  +\frac{h}{2}(n_y'^2+n_z'^2),
\end{equation}
式中，$f$，$g$，$h$ 是常数.

为了确定起见，我们现在来研究一个透镜，这时我们令 $n_x'\approx 1$；对于反射镜，以后可以知道，所有的公式都有相似的形式.现在将（56.1）代入通式（55.6）内，我们得到：
\begin{equation}
\begin{aligned}
  n_y(x-g)-fn_y'=y, &\qquad fn_y+n_y'(x'+h)=y',\\
  n_z(x-g)-fn_z'=z, &\qquad fn_z+n_z'(x'+h)=z'.
\end{aligned}
\end{equation}

我们考虑一个从点 $x,y,z$ 射出来的同心光束；设点 $x',y',z'$ 是光束的光线透过透镜后的交点.如果（56.2）的第一对方程与第二对方程都是独立的，那么，这四个方程当 $x,y,z,x',y',z'$ 为已知时决定了 $n_y,n_z,n_y',n_z'$ 的一组确定值，这意味着只有一条光线从点 $x,y,z$ 出发并且经过点 $x',y',z'$.因此，为了使所有从 $x,y,z$ 出发的光线都经过 $x',y',z'$，（56.2）就必须不是独立的，也就是说，这些方程中的一对可以从另一对推出来.这种相依性的必要条件显然是这些方程中，一对方程的系数与另一对方程的系数成比例.因此我们有
\begin{equation}
  \frac{x-g}{f}=-\frac{f}{x'+h}=\frac{y}{y'}=\frac{z}{z'};
\end{equation}
特别有
\begin{equation}
  (x-g)(x'+h)=-f^2.
\end{equation}

我们所求得的这些方程就是窄光束成像情形中，像坐标与物坐标所满足的关系.

光轴上的两点 $x=g$ 及 $x'=-h$ 称为光学系统的主焦点.我们来考虑平行于光轴的光线束.这种光的源点显然是光轴的无穷远点，即 $x=\infty$. 从（56.3）式看出，在这种情况下，$x'=-h$.因此，一个平行光线束经过光学系统后，会相交于主焦点.反之，一束从主焦点射出的光线，经过光学系统后就变为平行光了.

在方程（56.3）中，坐标 $x$ 和 $x'$ 都是从光轴上同一原点量起的.然而，如果选择相应的主焦点作为原点，并从这些不同的原点去测量物坐标及像坐标，会更加便利些.我们选择从相应的主焦点到光行进的一边的方向为正方向.用大写字母代表新的物坐标及像坐标，则我们有
\begin{equation*}
  X=x-g,\quad X'=x'+h,\quad Y=y,\quad Y'=y',\quad Z=z,\quad Z'=z'.
\end{equation*}

成像的方程（56.3）及（56.4）在新坐标中取下面的形式：
\begin{equation}
  XX'=-f^2,
\end{equation}
\begin{equation}
  \frac{Y'}{Y}=\frac{Z'}{Z}=\frac{f}{X}=-\frac{X'}{f}.
\end{equation}

$f$ 这个量称为这个系统的主焦距.

$Y'/Y$ 这个比值称为横向放大率.至于纵向放大率，因为坐标并不彼此成简单比例，它就必须写成微分形式，用来比较物元素的长（沿轴的方向）与相应的像元素的长.从（56.5）式，我们得到“纵向放大率”为
\begin{equation}
  \left|\frac{\mathrm{d}X'}{\mathrm{d}X}\right|
  =\frac{f^2}{X^2}=\left(\frac{Y'}{Y}\right)^2.
\end{equation}

由此可见，哪怕对于无限小的物体来说，也不可能得到几何相似的像.纵向放大率绝对不会等于横向放大率（除了全同成像这种平凡情形外）.

从光轴上 $X=f$ 点出发的光线束，又交于同一轴上 $X'=-f$ 点；这两个点称为主点.从方程（56.2）$(n_yX-fn_y'=Y,\ n_zX-fn_z'=Z)$ 显而易见，在这种情形 $(X=f,\ Y=Z=0)$ 下，我们得到方程 $n_y=n_y'$，$n_z=n_z'$.因此，每条从主点出发的光线，又在另一主点与光轴相交，其方向与原方向平行.

如果物和它的像的坐标是从主点量起，而不是从主焦点量起，那么，对于这些坐标 $\xi$ 和 $\xi'$，我们有
\begin{equation*}
  \xi'=X'+f,\qquad \xi=X-f.
\end{equation*}

将上式代入（56.5），很容易得到成像方程如下：
\begin{equation}
  \frac{1}{\xi}-\frac{1}{\xi'}=-\frac{1}{f}.
\end{equation}

可以证明，对于薄光学系统（例如反射镜或薄透镜），两个主点差不多重合.在这种情形下，方程（56.8）就特别方便，因为在这个方程中，$\xi$ 和 $\xi'$ 实际上从同一点量起.

如果焦距为正，那么，位于焦点前面 $(X>0)$ 的物所生成的像是正立的 $(Y'/Y>0)$；这种光学系统称为会聚的.如果 $f<0$，那么，当 $X>0$ 时，我们得到 $Y'/Y<0$，这就意味着物所生成的像是倒立的；这种光学系统称为发散的.

还有一个成像的极限情形没有被包含在公式（56.8）内；在这种情形下，所有三个系数 $f$，$g$，$h$ 都是无穷大（也就是说，光学系统有无限大的焦距，而其主焦点位于无限远处）.在（56.4）式中取 $f$，$g$，$h$ 趋近于无穷大的极限，我们得到
\begin{equation*}
  x'=\frac{h}{g}x+\frac{f^2-gh}{g}.
\end{equation*}
因为我们只对物和像与光学系统之间的距离有限的情形感兴趣，$f$，$g$，$h$ 必须这样地趋近于无穷大，使比值 $h/g$，$(f^2-gh)/g$ 为有限.用 $\alpha^2$ 和 $\beta$ 分别代表这两个比值，则我们有 $x'=\alpha^2x+\beta$.

对于其他两个坐标，我们现在从一般方程（56.7）得到：
\begin{equation*}
  \frac{y'}{y}=\frac{z'}{z}=\pm\alpha.
\end{equation*}
最后，再次从不同的原点测量坐标 $x$ 和 $x'$，即从轴上的某一任意点及这个点的像测量坐标 $x$ 和 $x'$，我们最终得到成像方程的如下简单形式：
\begin{equation}
  X'=\alpha^2X,\qquad Y'=\pm\alpha Y,\qquad Z'=\pm\alpha Z.
\end{equation}

因此，纵向放大率和横向放大率都是常数（然而像在几何上并不与物相似，因为两个放大率不相等）.这种成像的情形称为望远镜式.

我们为透镜求出的从（56.5）到（56.9）的所有方程，对于反射镜同样可以应用，甚至于对非轴对称的光学系统也可以应用，只要用来成像的窄光线束是在光轴的附近行进就可以了.同时，物和像的坐标 $x$ 的度量必须总是从相应的点（主焦点及主点）顺着光线传播方向沿光轴进行.这时，必须记住，在不具备轴对称性的光学系统中，光学系统前后的光轴方向不在同一直线上.

\begin{xiti}

\noindent{\heiti 习题1}\quad 利用两个光轴重合的轴对称光学系统成像，试求其焦距.

解：设 $f_1$ 和 $f_2$ 是这两个系统的焦距.对于每一个系统，我们分别有
\begin{equation*}
  X_1X_1'=-f_1^2,\qquad X_2X_2'=-f_2^2.
\end{equation*}
因为第一个系统所产生的像是第二个系统的物，那么，用 $l$ 来代表第一个系统后主焦点与第二个系统前焦点的距离，我们就有 $X_2=X_1'-l$.通过 $X_1$ 来表示 $X_2'$，得到
\begin{equation*}
  X_2'=\frac{X_1f_2^2}{f_1^2+lX_1},
\end{equation*}
或
\begin{equation*}
  \left(X_1+\frac{f_1^2}{l}\right)\left(X_2'-\frac{f_2^2}{l}\right)
  =-\left(\frac{f_1f_2}{l}\right)^2,
\end{equation*}
由此可见，组合系统的主焦点位于点 $X_1=-f_1^2/l$，$X_2'=f_2^2/l$，而焦距则是
\begin{equation*}
  f=-\frac{f_1f_2}{l}
\end{equation*}
（为了选择这个式子的正负号，我们必须写出相应的横向放大率的方程）.

在 $l=0$ 的情形，焦距 $f=\infty$，这就意味着由组合系统得到望远镜式的成像.在这种情形下，我们有 $X_2'=X_1(f_2/f_1)^2$，就是说，通式（56.9）中的参数 $\alpha=f_2/f_1$.

\noindent{\heiti 习题2}\quad 求带电粒子“磁透镜”的焦距，该磁透镜由长为 $l$ 的一段纵向均匀磁场所产生（图8）\footnote{在忽略螺线管两端附近场的均匀性的扰动时，螺线管内部就有这样的场.}.

解：粒子在磁场中运动时它的动能是守恒的；因而对于约化作用量 $S_0(\boldsymbol{r})$（总作用量为 $S=-\mathcal{E}t+S_0$）的哈密顿-雅可比方程是
\begin{equation*}
  \left(\nabla S_0-\frac{e}{c}\boldsymbol{A}\right)^2=p^2,
\end{equation*}
式中
\begin{equation*}
  p^2=\frac{\mathcal{E}^2}{c^2}-m^2c^2=\mathrm{const}.
\end{equation*}

\par\nopagebreak\noindent\begin{minipage}{\linewidth}\centering\includegraphics[width=4.5cm]{fig8.pdf}\\[0.3em]{\small 图 8}\end{minipage}\par\nopagebreak

用均匀磁场矢势的公式（19.4），选择沿场的方向为 $x$ 轴，并将该轴看成轴对称光学系统的光轴，我们得到形如下式的哈密顿-雅可比方程：
\begin{equation*}
  \left(\frac{\partial S_0}{\partial x}\right)^2
  +\left(\frac{\partial S_0}{\partial r}\right)^2
  +\frac{e^2}{4c^2}H^2r^2=p^2,
  \tag{1}
\end{equation*}
式中 $r$ 是离 $x$ 轴的距离，$S_0$ 是 $x$ 和 $r$ 的函数.

因为窄束粒子靠近光轴行进，坐标 $r$ 很小，所以我们可以尝试把 $S_0$ 表为 $r$ 的幂级数.这个级数的头两项是
\begin{equation*}
  S_0=px+\frac{1}{2}\sigma(x)r^2,
  \tag{2}
\end{equation*}
式中 $\sigma(x)$ 满足方程
\begin{equation*}
  p\sigma'(x)+\sigma^2+\frac{e^2}{4c^2}H^2=0.
  \tag{3}
\end{equation*}

在透镜前方的区域 1，我们得到：
\begin{equation*}
  \sigma^{(1)}=\frac{p}{x-x_1},
\end{equation*}
式中 $x_1<0$ 是常数.这个解对应于自由粒子束，沿着直线从区域 1 内光轴上 $x=x_1$ 点出射.实际上，对于从 $x=x_1$ 点出射的具有动量 $p$ 的自由运动粒子，作用量函数是
\begin{equation*}
  S_0=p\sqrt{r^2+(x-x_1)^2}
  \approx p(x-x_1)+\frac{pr^2}{2(x-x_1)}.
\end{equation*}

类似地，在透镜后方区域 2 中，我们写出：
\begin{equation*}
  \sigma^{(2)}=\frac{p}{x-x_2},
\end{equation*}
式中常数 $x_2$ 是点 $x_1$ 的像的坐标.

在透镜里面的区域 3，可以用分离变量法得到方程（3）的解，我们有：
\begin{equation*}
  \sigma^{(3)}=\frac{eH}{2c}\cot\left(\frac{eH}{2cp}x+C\right),
\end{equation*}
式中 $C$ 是一个任意常数.

常数 $C$ 和 $x_2$（对于给定的 $x_1$）可以借助 $\sigma(x)$ 在 $x=0$ 和 $x=l$ 的连续性要求定出：
\begin{equation*}
  -\frac{p}{x_1}=\frac{eH}{2c}\cot C,\qquad
  \frac{p}{l-x_2}=\frac{eH}{2c}\cot\left(\frac{eH}{2cp}l+C\right).
\end{equation*}

从这些方程消去常数 $C$，我们得到
\begin{equation*}
  (x_1-g)(x_2+h)=-f^2,
\end{equation*}
式中\footnote{所给的 $f$ 值有正确的正负号，不过为证明这一点，需要作更多的研究.}
\begin{equation*}
  g=-\frac{2cp}{eH}\cot\frac{eHl}{2cp},\qquad h=g+l,
\end{equation*}
\begin{equation*}
  f=\frac{2cp}{eH\sin\dfrac{eHl}{2cp}}.
\end{equation*}

\end{xiti}

\section{宽光线束成像}

我们在上节所考虑的用窄光线束来成像的情形是近似的；光束愈窄，像愈准确（即像愈清晰）.我们现在来讨论利用任意宽的光线束给物体成像的问题.

用窄光线束给物体成像对于任何具有轴对称性的光学系统都可能办到，与此情形不同，利用宽光线束来成像仅仅对于特别构造的光学系统才有可能.在 §56 中已经指出过，即使加上了这个限制，成像也不是在空间所有的点都可能.

以后的推导都基于下面的重要说明.设所有从某一点出发的光线，穿过光学系统后，又在另外一点 $O'$ 相交.很容易看出，所有光线的光程长 $\psi$ 都是一样的.实际上，在 $O$ 或 $O'$ 每一点的近旁，相交的光线的波面是分别以 $O$ 及 $O'$ 为心的球面，且在趋近于 $O$ 和 $O'$ 的极限情形下，波面退化为这两个点.但是波面是等相位面，因此，沿着不同的光线，在它们与两个已定波面的交点之间的相位变化是一样的.从上面所说的可以推断，不同的光线在 $O$ 和 $O'$ 两点间的总的相位变化也是一样的.

首先让我们来考虑利用宽光线束使一小段直线生成像所必须满足的条件；这时，像也是一小段直线.我们选择这两个线段的方向为 $\xi$ 和 $\xi'$ 轴的方向，其原点相应地在物点 $O$ 和像点 $O'$ 上. 设 $\psi$ 为从 $O$ 出发到达 $O'$ 的光线的光程长度.如果一条光线从与 $O$ 点无限近的一点（其坐标为 $\mathrm{d}\xi$）出发，到达像所在的一点（其坐标为 $\mathrm{d}\xi'$），这条光线的光程长将是 $\psi+\mathrm{d}\psi$，其中
\begin{equation*}
  \mathrm{d}\psi=\frac{\partial\psi}{\partial\xi}\mathrm{d}\xi
  +\frac{\partial\psi}{\partial\xi'}\mathrm{d}\xi'.
\end{equation*}

我们引入“放大率”
\begin{equation*}
  \alpha_\xi=\frac{\mathrm{d}\xi'}{\mathrm{d}\xi}
\end{equation*}
它是像元的长 $\mathrm{d}\xi'$ 与物元的长 $\mathrm{d}\xi$ 之比.因为物的线段很短，放大率 $\alpha$ 沿着该线段可以当做是常数.按照平常的写法，$\partial\psi/\partial\xi=-n_\xi$，$\partial\psi/\partial\xi'=n_\xi'$（$n_\xi$，$n_\xi'$ 是光线与相应的 $\xi$ 轴和 $\xi'$ 轴夹角的余弦），我们得到
\begin{equation*}
  \mathrm{d}\psi=(\alpha_\xi n_\xi'-n_\xi)\,\mathrm{d}\xi.
\end{equation*}

对于每一对物与像的相应点，光程长 $\psi+\mathrm{d}\psi$ 对于所有从点 $\mathrm{d}\xi$ 出发而到达点 $\mathrm{d}\xi'$ 的光线都是一样的.从此我们得到一个条件：
\begin{equation}
  \alpha_\xi n_\xi'-n_\xi=\mathrm{const}.
\end{equation}

这就是我们所求的条件，它是在利用宽光线束使一段直线成像时，光线在光学系统内的路程所必须满足的条件.所有从 $O$ 点出发的光线都必须满足（57.1）式.

让我们应用条件（57.1）到利用轴对称光学系统成像的情形.

从与系统的光轴（$x$ 轴）重合的线段成像开始，显而易见，像也与光轴重合.由于光学系统的轴对称性，一条沿光轴行进的光线 $(n_x=1)$ 通过光学系统后并不改变它的方向，也就是说，$n_x'=1$.由此可以断定，（57.1）式中的常数在这种情况下等于 $\alpha_x-1$；我们可以将（57.1）式改写为
\begin{equation*}
  \frac{1-n_x}{1-n_x'}=\alpha_x.
\end{equation*}

用 $\theta$ 和 $\theta'$ 代表光线在物与像的所在点与光轴的夹角，则我们有
\begin{equation*}
  1-n_x=1-\cos\theta=2\sin^2\frac{\theta}{2},\qquad
  1-n_x'=1-\cos\theta'=2\sin^2\frac{\theta'}{2}.
\end{equation*}

因此，我们得到成像的条件如下：
\begin{equation}
  \frac{\sin\dfrac{\theta}{2}}{\sin\dfrac{\theta'}{2}}
  =\mathrm{const}=\sqrt{\alpha_x}.
\end{equation}

接下来，我们来研究垂直于轴对称系统光轴的一小块平面的成像；该像显然也垂直于这个轴.将（57.1）用于待成像平面内的任意线段.再用 $\theta$ 和 $\theta'$ 代表光线与光轴的夹角，我们得到：
\begin{equation*}
  \alpha_r\sin\theta'-\sin\theta=\mathrm{const},
\end{equation*}
对于从物平面与光轴的交点出射而又沿着这个光轴行进的光线 $(\theta=0)$，根据对称性，我们必定有 $\theta'=0$. 因此，$\mathrm{const}=0$，于是我们得到成像条件为
\begin{equation}
  \frac{\sin\theta}{\sin\theta'}=\mathrm{const}=\alpha_r.
\end{equation}

至于利用宽光线束来形成三维物体的像，很容易看出，即使对于一个无限小的体积，这也是不可能的，因为条件（57.2）和（57.3）不能兼容.

\section{几何光学的极限}

直接从单色平面波的定义可以知道，这种波的振幅无论在任何位置和任何时刻都是一样的.这样的波在空间的任何方向都无限延伸，而且在从 $-\infty$ 到 $+\infty$ 的整个时间范围内都存在.假如波的振幅在所有位置和所有时刻不保持为常量，那么，这种波最多只能是近乎于单色的.我们现在来讨论一个波的“非单色程度”的问题.

让我们来考虑一个振幅为时间函数的电磁波.换句话说，在波经过的空间每一点，波的振幅随着时间变化.设 $\omega_0$ 是波的某一平均频率.这时，波的场，例如电场，在一指定点有 $E_0(t)\mathrm{e}^{-\mathrm{i}\omega_0 t}$ 的形式.这个场本身当然不是单色的，然而可以展开为单色波，即可以展开为傅里叶积分.这个展开式中，分量的振幅与积分
\begin{equation*}
  \int_{-\infty}^{+\infty}E_0(t)\mathrm{e}^{\mathrm{i}(\omega-\omega_0)t}\,\mathrm{d}t
\end{equation*}
成正比，$\omega$ 是展开式的分量的频率.因子 $\mathrm{e}^{\mathrm{i}(\omega-\omega_0)t}$ 是一个周期性函数，其平均值为零. 如果 $E_0$ 是常量，那么，对于所有 $\omega\neq\omega_0$ 的角频率，积分就真正等于零了.如果 $E_0(t)$ 是变量，但是在一个大小与 $1/|\omega-\omega_0|$ 同量级的时间间隔内变化很小，那么，这个积分就几乎等于零，$E_0$ 的变化愈慢，积分就愈近于零.为了使积分与零有显著的不同，$E_0(t)$ 在一个与 $1/|\omega-\omega_0|$ 同量级的时间间隔内必须有显著的变化.

我们用 $\Delta t$ 表示这样一个时间间隔的量级，在这个时间间隔内，波的振幅在空间一给定点的变化是显著的.从这些考虑，现在可以推断，那些在这个波的谱分解中有显著强度的、与 $\omega_0$ 相差最大的频率将由 $1/|\omega-\omega_0|\sim\Delta t$ 这个条件决定.如果用 $\Delta\omega$ 表示平均频率 $\omega_0$ 附近的频率间隔，而这个频率间隔是在波的谱分解之内，那么，我们将有关系式
\begin{equation}
  \Delta\omega\cdot\Delta t\sim 1.
\end{equation}

由此可见，$\Delta t$ 愈大，波愈近于单色（亦即 $\Delta\omega$ 愈小），也就是说，波在空间给定点的振幅变化也就愈慢.

对于波矢也容易推导出与（58.1）相似的关系.设 $\Delta x$，$\Delta y$，$\Delta z$ 是沿着 $x,y,z$ 各轴的距离的数量级，在这些距离内，波的振幅有显著的变化.在给定时刻，波的场作为坐标 $x$ 的函数（$y$ 和 $z$ 固定）有如下形式
\begin{equation*}
  E_0(\boldsymbol{r})\,\mathrm{e}^{\mathrm{i}\boldsymbol{k}_0\cdot\boldsymbol{r}},
\end{equation*}
此处的 $\boldsymbol{k}_0$ 是波矢的某个平均值.与推导（58.1）式完全一样，我们可以求得波的傅里叶积分展开式中所含的间隔值 $\Delta\boldsymbol{k}$.这时，我们得到
\begin{equation}
  \Delta k_x\cdot\Delta x\sim 1,\qquad
  \Delta k_y\cdot\Delta y\sim 1,\qquad
  \Delta k_z\cdot\Delta z\sim 1.
\end{equation}

让我们来研究在有限时间间隔内辐射出去的波这个特例.用 $\Delta t$ 代表这个时间间隔的数量级.波在空间给定点的振幅在波完全通过该点的时段 $\Delta t$ 内有显著的变化.根据（58.1）式，我们现在可以说这样的波的“非单色性” $\Delta\omega$ 不能小于 $1/\Delta t$（大一些当然是可以的）：
\begin{equation}
  \Delta\omega\gtrsim\frac{1}{\Delta t}.
\end{equation}

同样，如果 $\Delta x$，$\Delta y$，$\Delta z$ 是波在空间延展的数量级，那么，对于在波的分解中的波矢分量值的分布，我们得到
\begin{equation}
  \Delta k_x\gtrsim\frac{1}{\Delta x},\qquad
  \Delta k_y\gtrsim\frac{1}{\Delta y},\qquad
  \Delta k_z\gtrsim\frac{1}{\Delta z}.
\end{equation}

从这些公式可以断定，如果光束有有限的宽度，那么，在这样的光束中，光的传播方向不可能严格不变.取光束中光的（平均）方向作为 $x$ 轴，我们得到
\begin{equation}
  \theta_y\gtrsim\frac{1}{k\Delta y}\sim\frac{\lambda}{\Delta y},
\end{equation}
式中，$\theta_y$ 是光束在 $xy$ 平面内对其平均方向偏差的数量级，$\lambda$ 是波长.

另一方面，公式（58.5）回答了光学成像清晰度的极限问题.按照几何光学，一束光束的所有光线都应相交在一点，而在实际上，所得的像并不是一个几何点，而是形成一个斑点.按照（58.5）式，我们得到这个斑点的宽度 $\Delta$
\begin{equation}
  \Delta\sim\frac{1}{k\theta}\sim\frac{\lambda}{\theta},
\end{equation}
其中，$\theta$ 是光束的张角.这些公式不但可以应用到像上，而且也可以应用到物上.就是说，我们可以断定，在观察从一个发光点出射的光束时，不可能区别这个点与一个大小为 $\lambda/\theta$ 的物体.相应地，公式（58.6）决定显微镜的分辨能力的极限.$\Delta$ 的最小值是 $\lambda$（在 $\theta\sim 1$ 时达到），这与如下事实完全符合，即几何光学的极限是由光波的波长决定的.

\begin{xiti}

\noindent{\heiti 习\ 题}\quad 求一个与光阑相距 $l$ 的平行光束所产生的光束的最小宽度的数量级.

解：设在光阑上的孔径为 $d$，从（58.5）得到光束的偏转角（“衍射角”）为 $\lambda/d$，从而得到光束宽度的数量级为 $d+(\lambda/d)l$.这个量的最小值 $\sim\sqrt{\lambda l}$.

\end{xiti}

\section{衍射}

几何光学定律只有在波长可以看做无限小的理想情况下才是严格正确的.这个条件满足得愈不好，与几何光学的偏差就愈大.由这种差异所导致的现象被称为衍射现象.

例如，在光\footnote{以后我们所讨论的衍射，都是光的衍射；所有这些同样的讨论，当然也可以应用于任意的电磁波.}传播的路径上，放置一个障碍物——一个任意形状的不透明物体（我们称它为屏），或者，例如，光经过不透明屏的孔，这时，我们就能观察到衍射现象.如果几何光学定律被严格地满足的话，那么，在屏后的影区与光照区域会有非常清楚的分界线.由于衍射，光与影之间没有明晰的界限，而光强分布图形甚为复杂.屏愈小，或屏上的光孔愈小，或波长愈大，这种衍射现象就愈强.

衍射理论的任务就是，在物体位置与形状一定（光源位置也一定）的情况下，求光的分布，亦即整个空间电磁场的分布.这个问题的严格解只有通过求解波动方程才可能得到，求解时要利用到物体表面所满足的适当边界条件，而这些边界条件则与物质的光学性质有关.这样的求解往往遇到很大的数学困难.

然而，在大多数情况下，对于光在光与影的界限附近的分布问题，利用近似方法求解就够了.我们可以将这种方法用于与几何光学相差不大的情形，这种情形是：第一，所有物体的大小都比波长大很多（这个要求既适用于屏和孔的尺度，也适用于从物体到光的发射点和观察点的距离）；第二，光的方向与几何光学所给出的光线方向相差很小.

现在我们来考虑任意一个有孔的屏，光线从光源出发穿过这个孔.图9表示屏的断面（粗线）；光行进的方向是从左向右.用 $u$ 代表 $\boldsymbol{E}$ 或 $\boldsymbol{H}$ 的任一分量.这时，我们把 $u$ 理解为仅仅是坐标的函数，即不包含与时间有关的因子 $\mathrm{e}^{-\mathrm{i}\omega t}$.我们的问题是决定光的强度，也就是决定在屏后面的任何一个观察点 $P$ 的场 $u$.在与几何光学相差无几的情况下近似求解这个问题时，我们可以认为，在屏孔的点上，场与没有屏存在时一样.换句话说，此处的场可以直接从几何光学推出.屏背面的所有点上，场可以算做零.这时，屏本身的特性（即屏物质的特性）显然不起作用.同时也很明显，在我们所考虑的情形下，对于衍射，重要的仅是屏孔的边缘形状，至于不透明的屏的形状则无关紧要.

我们用任一曲面盖着屏的孔，而这个曲面以孔的边缘为界（这样的一个曲面的断面图用虚线表示在图9中）.我们将这个曲面分割为面积为 $\mathrm{d}f$ 的若干块，$\mathrm{d}f$ 的尺度比孔小，但比光的波长大.我们可以将这些小块（光要经过这些小块）中的每一块本身当做光波的源，光波从这个小块向各方向射出.我们来考虑在 $P$ 点的场，这一点的场是遮盖着屏孔表面的所有小块 $\mathrm{d}f$ 在该点所产生的场的叠加的结果（这被称为惠更斯原理）.

\par\nopagebreak\noindent\begin{minipage}{\linewidth}\centering\includegraphics[width=4.5cm]{fig9.pdf}\\[0.3em]{\small 图 9}\end{minipage}\par\nopagebreak

小块 $\mathrm{d}f$ 在 $P$ 点所产生的场显然与在小块 $\mathrm{d}f$ 本身上的场之值 $u$ 成正比（提醒一下，我们已经假定了，在 $\mathrm{d}f$ 本身上的场与没有屏存在时一样）.此外，在 $P$ 点的场又与面积 $\mathrm{d}f$ 在与方向 $\boldsymbol{n}$ 垂直的平面上的投影 $\mathrm{d}f_n$ 成正比，方向 $\boldsymbol{n}$ 是光从光源到 $\mathrm{d}f$ 的方向.这是从以下的事实得来的，即不管面元 $\mathrm{d}f$ 的形状如何，只要投影面 $\mathrm{d}f_n$ 保持不变，相同的光线将经过面元 $\mathrm{d}f$；因此面元 $\mathrm{d}f$ 对 $P$ 点的场的影响将是一样的.

由此可见，小块 $\mathrm{d}f$ 在 $P$ 点所产生的场与 $u\,\mathrm{d}f_n$ 成正比.此外，我们还必须考虑到波从 $\mathrm{d}f$ 传播到 $P$ 点的过程中的相位与振幅的变化.这种变化的规律是由公式（54.3）来决定的.因此必须以 $(1/R)\mathrm{e}^{\mathrm{i}kR}$ 乘 $u\,\mathrm{d}f_n$（此处的 $R$ 是从 $\mathrm{d}f$ 到 $P$ 的距离，而 $k$ 则是光的波矢的绝对值），于是我们得到所求的场等于
\begin{equation*}
  au\,\frac{\mathrm{e}^{\mathrm{i}kR}}{R}\,\mathrm{d}f_n,
\end{equation*}
其中 $a$ 仍是一个未知常数.$P$ 点的场是所有面元 $\mathrm{d}f$ 在该点所产生的场的叠加的结果，因此它等于
\begin{equation}
  u_P=a\int\frac{u\,\mathrm{e}^{\mathrm{i}kR}}{R}\,\mathrm{d}f_n,
\end{equation}
此处的积分遍及以屏孔边缘为界的曲面.在我们所考虑的近似中，这个积分当然不能与这个曲面的形状有关.公式（59.1）显然不仅可以应用到屏上有孔的衍射，而且也可以应用到光在其周围可以自由传播的屏的衍射.在这种情形下，（59.1）式内的积分面应该伸延到屏的各个边线.

为了决定常数 $a$，我们考虑一个沿 $x$ 轴传播的平面波；波面与 $yz$ 平面平行.设 $u$ 为场在 $yz$ 平面内的值.那么，$P$ 点（我们选择这一点在 $x$ 轴上）的场等于 $u_P=u\,\mathrm{e}^{\mathrm{i}kx}$.另一方面，$P$ 点的场也可以由公式（59.1）来决定，选取（例如）$yz$ 平面为积分面.同时，因为衍射角很小，在 $yz$ 平面内的点中，只有那些靠近原点的点在积分中才重要，即重要的是 $y,z\ll x$ 的点（$x$ 是 $P$ 点的坐标）.因此，
\begin{equation*}
  R=\sqrt{x^2+y^2+z^2}\approx x+\frac{y^2+z^2}{2x}
\end{equation*}
而（59.1）则化为
\begin{equation*}
  u_P=au\,\frac{\mathrm{e}^{\mathrm{i}kx}}{x}
  \int_{-\infty}^{+\infty}\exp\left(\mathrm{i}k\frac{y^2}{2x}\right)\mathrm{d}y
  \cdot\int_{-\infty}^{+\infty}\exp\left(\mathrm{i}k\frac{z^2}{2x}\right)\mathrm{d}z,
\end{equation*}
式中，$u$ 是常数（在 $yz$ 平面内的场）；在因子 $1/R$ 中，我们可以使 $R\approx x=\mathrm{const}$. 若将 $y=\xi\sqrt{2x/k}$ 代入，则每个积分都变为
\begin{equation*}
  \int_{-\infty}^{+\infty}\mathrm{e}^{\mathrm{i}\xi^2}\mathrm{d}\xi
  =\int_{-\infty}^{+\infty}\cos\xi^2\,\mathrm{d}\xi
  +\mathrm{i}\int_{-\infty}^{+\infty}\sin\xi^2\,\mathrm{d}\xi
  =\sqrt{\frac{\pi}{2}}\,(1+\mathrm{i}),
\end{equation*}
从而我们得到 $u_P=au\,\mathrm{e}^{\mathrm{i}kx}\cdot 2\mathrm{i}\pi/k$.另一方面，$u_P=u\,\mathrm{e}^{\mathrm{i}kx}$，因此 $a=k/(2\pi\mathrm{i})$. 将它代入（59.1）式，我们最后得到所提出问题的解如下：
\begin{equation}
  u_P=\int\frac{ku}{2\pi\mathrm{i}R}\,\mathrm{e}^{\mathrm{i}kR}\,\mathrm{d}f_n.
\end{equation}

在推导（59.2）式时，假设了光源实质上是一个点，并且假设了光是严格单色的。然而，对于发射非单色光的实际延展光源的情形，并不需要做特殊的处理。由于光源不同点发射的光完全独立（非相干），以及发射光不同谱成分的非相干性，总的衍射图案就是从光的各独立成分衍射得到的强度分布之和.

现在让我们应用公式（59.2）来求一条光线在经过它与焦散面相切之点时的相位变化问题（见 §54 末尾）.我们选择任意的一个波面作为在（59.2）式中的积分面，再求 $P$ 点的场 $u_P$，此处的 $P$ 点是在某一条光线上，它与光线同我们所选定的波面的交点的距离为 $x$（我们选定这一点作为坐标原点 $O$，而以在 $O$ 点与波面相切的平面作为 $yz$ 平面）.在求（59.2）的积分时，仅仅波面在 $O$ 点附近的一个小面积是重要的.如果选择 $xy$ 和 $xz$ 平面与波面在 $O$ 点的主曲率平面相重合，那么，在这一点的附近，波面的方程是
\begin{equation*}
  X=\frac{y^2}{2R_1}+\frac{z^2}{2R_2},
\end{equation*}
式中，$R_1$ 和 $R_2$ 是曲率半径.从在波面上以 $X,y,z$ 为坐标的一点到以 $x,0,0$ 为坐标的 $P$ 点的距离 $R$ 则是
\begin{equation*}
  R=\sqrt{(x-X)^2+y^2+z^2}
  \approx x+\frac{y^2}{2}\left(\frac{1}{x}-\frac{1}{R_1}\right)
  +\frac{z^2}{2}\left(\frac{1}{x}-\frac{1}{R_2}\right).
\end{equation*}

在波面上，场 $u$ 可以当做常数；因子 $1/R$ 亦可当做常数.既然我们只对波的相位变化有兴趣，所以我们略去系数，并简单地写
\begin{equation}
\begin{aligned}
  u_P\sim\frac{1}{\mathrm{i}}\int\mathrm{e}^{\mathrm{i}kR}\,\mathrm{d}f_n
  &\approx\frac{\mathrm{e}^{\mathrm{i}kx}}{\mathrm{i}}
  \int_{-\infty}^{+\infty}\exp\left[\mathrm{i}k\frac{y^2}{2}
  \left(\frac{1}{x}-\frac{1}{R_1}\right)\right]\mathrm{d}y\\
  &\qquad\cdot\int_{-\infty}^{+\infty}\exp\left[\mathrm{i}k\frac{z^2}{2}
  \left(\frac{1}{x}-\frac{1}{R_2}\right)\right]\mathrm{d}z.
\end{aligned}
\end{equation}

波面的曲率中心是在我们所考虑的光线上的 $x=R_1$ 和 $x=R_2$ 两点；这两点就是光线与焦散面相切的切点.设 $R_2<R_1$.当 $x<R_2$ 时，两个积分号内的指数式中的 $\mathrm{i}$ 的系数都是正数，这两个积分的每一个都正比于 $(1+\mathrm{i})$.因此，在未与焦散面第一次相切的这部分光线上，我们有 $u_P\sim\mathrm{e}^{\mathrm{i}kx}$.当 $R_2<x<R_1$ 时，在两个切点之间的这部分光线上，沿着 $y$ 的积分正比于 $(1+\mathrm{i})$，但沿着 $z$ 的积分则正比于 $1-\mathrm{i}$，因此两者之积就不包含 $\mathrm{i}$. 于是在这里我们有 $u_P\sim-\mathrm{i}\,\mathrm{e}^{\mathrm{i}kx}=\mathrm{e}^{\mathrm{i}(kx-\pi/2)}$，就是说，当光线经过第一个焦散面附近时，光线的相位发生了 $-\pi/2$ 的额外变化.最后，当 $x>R_1$ 时，我们有 $u_P\sim-\mathrm{e}^{\mathrm{i}kx}=\mathrm{e}^{\mathrm{i}(kx-\pi)}$，就是说当光线经过第二个焦散面附近时，相位又发生了一次 $-\pi/2$ 的变化.

\begin{xiti}

\noindent{\heiti 习\ 题}\quad 求光线与焦散面相切的切点附近的光强度分布.

解：为了解决这个问题，我们利用公式（59.2），其中的积分面可以用任意的一个波面，但这个波面离开光线同焦散面的切点要足够得远.在图10中，$ab$ 是这个波面的一段，而 $a'b'$ 是焦散面的一段；$a'b'$ 是曲线 $ab$ 的渐屈线.我们所要研究的是光线 $QO$ 与焦散面的切点 $O$ 附近的光强度分布；假设光线 $QO$ 段的长度 $D$ 足够大.我们用 $x$ 代表从 $O$ 点沿着焦散面的法线的距离，并认为在法线上的点的 $x$ 值向着曲率中心的方向为正.

\par\nopagebreak\noindent\begin{minipage}{\linewidth}\centering\includegraphics[width=6.5cm]{fig10.pdf}\\[0.3em]{\small 图 10}\end{minipage}\par\nopagebreak

（59.2）式中的被积函数是距离 $R$ 的函数，$R$ 是从波面上任意一点 $Q'$ 到 $P$ 点的距离.按照熟知的渐屈线的性质，弧 $OO'$ 的长与线段 $Q'O'$ 的长之和（$Q'O'$ 切于 $O'$ 点）等于线段 $QO$ 之长（$QO$ 切于 $O$ 点）.对于互相邻近的点 $O$ 与 $O'$，我们有 $OO'=\theta\rho$（$\rho$ 是焦散面在 $O$ 点的曲率半径）. 因此 $Q'O'=D-\theta\rho$.距离 $Q'O$（沿直线）近似地等于（假设 $\theta$ 角很小）：
\begin{equation*}
  Q'O\approx Q'O'+\rho\sin\theta
  =D-\theta\rho+\rho\sin\theta\approx D-\frac{\rho\theta^3}{6}.
\end{equation*}
最后，距离 $R=Q'P$ 是 $R\approx Q'O-x\sin\theta\approx Q'O-x\theta$，即
\begin{equation*}
  R\approx D-x\theta-\frac{1}{6}\rho\theta^3.
\end{equation*}

将上式代入（59.2），我们得到
\begin{equation*}
  u_P\sim\int_{-\infty}^{+\infty}
  \exp\left(-\mathrm{i}kx\theta-\frac{\mathrm{i}k\rho}{6}\theta^3\right)\mathrm{d}\theta
  =2\int_0^{\infty}\cos\left(kx\theta+\frac{k\rho}{6}\theta^3\right)\mathrm{d}\theta
\end{equation*}
（被积函数内的 $1/D$ 的变化是很慢的，因而比起指数因子来就不重要了，所以我们假设它是常数）.引入新积分变量 $\xi=(k\rho/2)^{1/3}\theta$，我们得到
\begin{equation*}
  u_P\sim\Phi\left(x\sqrt[3]{\frac{2k^2}{\rho}}\right),
\end{equation*}
式中，$\Phi(t)$ 即所谓艾里函数\footnote{艾里函数 $\Phi(t)$ 的定义是
\begin{equation*}\tag{1}
  \Phi(t)=\frac{1}{\sqrt{\pi}}\int_0^{\infty}\cos\left(\frac{\xi^3}{3}+\xi t\right)\mathrm{d}\xi
\end{equation*}
（见本教程第三卷 §b）.当函数的自变量是大的正值时，$\Phi(t)$ 的渐近式是
\begin{equation*}\tag{2}
  \Phi(t)\approx\frac{1}{2t^{1/4}}\exp\left(-\frac{2}{3}t^{3/2}\right).
\end{equation*}
就是说 $\Phi(t)$ 按照指数式趋近于零.当 $t$ 是大的负值时，$\Phi(t)$ 按规律
\begin{equation*}\tag{3}
  \Phi(t)\approx\frac{1}{(-t)^{1/4}}\sin\left[\frac{2}{3}(-t)^{3/2}+\frac{\pi}{4}\right].
\end{equation*}
作减幅振荡.
艾里函数与 $1/3$ 阶麦克唐纳函数（变型汉克尔函数）的关系为：
\begin{equation*}\tag{4}
  \Phi(t)=\sqrt{\frac{t}{3\pi}}\,K_{1/3}\left(\frac{2}{3}t^{3/2}\right).
\end{equation*}
公式（2）对应于 $K_\nu(t)$ 的渐近展开：
\begin{equation*}
  K_\nu(t)\approx\sqrt{\frac{\pi}{2t}}\,\mathrm{e}^{-t}.
\end{equation*}}.

% ---- 脚注①（原书 p.168）已由脚注代理对照原书扫描页转录排入（2026-10-10） ----

对于强度 $I\sim|u_P|^2$，我们可写出
\begin{equation*}
  I=2A\left(\frac{2k^2}{\rho}\right)^{1/6}
  \Phi^2\left(x\sqrt[3]{\frac{2k^2}{\rho}}\right)
\end{equation*}
（关于常数因子的选择，见下）.

当 $x$ 是大的正值时，我们由此得到渐近公式
\begin{equation*}
  I\approx\frac{A}{2\sqrt{x}}
  \exp\left(-\frac{4x^{3/2}}{3}\sqrt{\frac{2k^2}{\rho}}\right),
\end{equation*}
就是说强度按指数规律下降（影区）.当 $x$ 是大的负值时，我们有
\begin{equation*}
  I\approx\frac{2A}{\sqrt{-x}}
  \sin^2\left[\frac{2(-x)^{3/2}}{3}\sqrt{\frac{2k^2}{\rho}}
  +\frac{\pi}{4}\right],
\end{equation*}
就是说强度快速地振荡；强度在振荡时的平均值是
\begin{equation*}
  \overline{I}=\frac{A}{\sqrt{-x}}.
\end{equation*}

由此显见常数 $A$ 的意义——它是忽略衍射效应时从几何光学得到的远离焦散面的强度.

函数 $\Phi(t)$ 在 $t=-1.02$ 达到其最大值 $0.949$，因而在 $x(2k^2/\rho)^{1/3}=-1.02$ 达到最大强度
\begin{equation*}
  I=2.03\,A\,k^{1/3}\rho^{-1/6}.
\end{equation*}

在光线与焦散面相切那一点 $(x=0)$，我们有
\begin{equation*}
  I=0.89\,A\,k^{1/3}\rho^{-1/6}
\end{equation*}
（因为 $\Phi(0)=0.629$）.所以在焦散面附近，光强度正比于 $k^{1/3}$，即正比于 $\lambda^{-1/3}$（$\lambda$ 是波长）.当 $\lambda\to 0$ 时，光强度趋向无穷大，正如它所应该的那样（见 §54）.

\end{xiti}

\section{菲涅耳衍射}

如果光源和我们求光强度的那一点 $P$ 都位于与屏的距离为有限之处，那么，在求 $P$ 点光的强度时，我们在（59.2）中取积分的波面上，只有一个小区域内的点才是重要的，这个区域是在光源和 $P$ 点的连线附近.实际上，既然同几何光学的偏差不大，从波面的各点到达 $P$ 点的光的强度，当我们从这条线移开时，就减小得很快.这种只有波面的一小部分起作用的衍射现象称为菲涅耳衍射现象.

现在我们来考虑任意一个屏的菲涅耳衍射.上面我们已经说过，对于给定点 $P$，只有在屏的边缘的一个小区域对于这种衍射是重要的.但是，对于足够小的区域，屏的边缘总可以当做一条直线.从现在起，凡说到屏的边缘意思就是指这样一小段直线.

我们选择经过光源 $Q$ 和屏的边缘线的平面作为 $xy$ 平面（图11）.我们选择 $xz$ 平面同 $xy$ 平面垂直并且经过 $Q$ 点和观察点 $P$，$P$ 点就是我们求光强度的那点.最后我们选择坐标原点 $O$ 在屏的边缘上，这样一来，所有三个轴的位置就完全确定了.

\par\nopagebreak\noindent\begin{minipage}{\linewidth}\centering\includegraphics[width=7cm]{fig11.pdf}\\[0.3em]{\small 图 11}\end{minipage}\par\nopagebreak

设从光源 $Q$ 到原点的距离为 $D_q$. 我们用 $D_p$ 代表观察点 $P$ 的 $x$ 坐标，而用 $d$ 代表 $P$ 点的 $z$ 坐标，即从 $P$ 点到 $xy$ 平面的距离.按照几何光学，光只能穿过 $xy$ 平面上方的点；至于在 $xy$ 平面下方的区域，按照几何光学应当处于阴影中（几何影区）.

现在我们来确定几何影区边缘附近屏上的光强度分布，即是在 $d$ 的值比 $D_p$ 和 $D_q$ 都小之处的光强度分布.负的 $d$ 值表示 $P$ 点位于几何影区中.

我们选择经过屏边缘线且垂直于 $xy$ 平面的半平面为（59.2）中的积分面.在这个面上的点的 $x$ 和 $y$ 坐标有 $x=y\tan\alpha$ 的关系（$\alpha$ 是屏边缘线与 $y$ 轴的夹角），而 $z$ 坐标是正的.光源 $Q$ 在相距为 $R_q$（从 $Q$ 所在点算起）的地方所产生的波的场正比于因子 $\exp(\mathrm{i}kR_q)$.因此在积分面上的场 $u$ 正比于
\begin{equation*}
  u\sim\exp\left[\mathrm{i}k\sqrt{y^2+z^2+(D_q+y\tan\alpha)^2}\right].
\end{equation*}

在积分式（59.2）中，现在 $R$ 必须代以
\begin{equation*}
  R=\sqrt{y^2+(z-d)^2+(D_p-y\tan\alpha)^2}.
\end{equation*}

被积函数中的缓变因子比起指数因子来是不重要的.因此，我们可以将 $1/R$ 当做常数，而以 $\mathrm{d}y\mathrm{d}z$ 代替 $\mathrm{d}f_n$. 于是我们求得在 $P$ 点的场为
\begin{equation}
\begin{aligned}
  u_P\sim\int_{-\infty}^{+\infty}\int_0^{\infty}
  \exp\Bigl[\mathrm{i}k\Bigl(\sqrt{(D_q+y\tan\alpha)^2+y^2+z^2}+\\
  +\sqrt{(D_p-y\tan\alpha)^2+y^2+(z-d)^2}\Bigr)\Bigr]\mathrm{d}y\,\mathrm{d}z.
\end{aligned}
\end{equation}

我们已经说过，落到 $P$ 点的光线主要是来自积分平面上 $O$ 点附近的点.因此在积分（60.1）中，只有那些比 $D_q$ 和 $D_p$ 都小的 $y$ 和 $z$ 的值才是重要的.

根据这个理由，可以写出
\begin{equation*}
\begin{aligned}
  \sqrt{(D_q+y\tan\alpha)^2+y^2+z^2}
  &\approx D_q+\frac{y^2+z^2}{2D_q}+y\tan\alpha,\\
  \sqrt{(D_p-y\tan\alpha)^2+y^2+(z-d)^2}
  &\approx D_p+\frac{y^2+(z-d)^2}{2D_p}-y\tan\alpha.
\end{aligned}
\end{equation*}

我们将上式代入（60.1）.因为我们只注意作为距离 $d$ 的函数的场，所以常数因子 $\exp[\mathrm{i}k(D_p+D_q)]$ 可以略去；对 $\mathrm{d}y$ 的积分所得的式子也不包含 $d$，因而也可以略去，于是我们有
\begin{equation*}
  u_P\sim\int_0^{\infty}\exp\left[
  \mathrm{i}k\left(\frac{1}{2D_q}z^2+\frac{1}{2D_p}(z-d)^2\right)\right]\mathrm{d}z.
\end{equation*}

这个式子也可以写成
\begin{equation}
  u_P\sim\exp\left\{\mathrm{i}k\frac{d^2}{2(D_p+D_q)}\right\}
  \int_0^{\infty}\exp\left\{\mathrm{i}k\,
  \frac{\dfrac{1}{2}\left[\left(\dfrac{1}{D_p}+\dfrac{1}{D_q}\right)z-\dfrac{d}{D_p}\right]^2}
       {\dfrac{1}{D_p}+\dfrac{1}{D_q}}\right\}\mathrm{d}z.
\end{equation}

光的强度由场的平方所决定，就是说，由 $u_P$ 的模的平方 $|u_P|^2$ 决定. 因此，当计算强度时，积分号前面的因子不存在了，因为用它的复共轭式来乘，其结果为1.在积分内，我们作以下的替换：
\begin{equation*}
  \frac{k}{2}\,
  \frac{\left[\left(\dfrac{1}{D_p}+\dfrac{1}{D_q}\right)z-\dfrac{d}{D_p}\right]^2}
       {\dfrac{1}{D_p}+\dfrac{1}{D_q}}=\eta^2
\end{equation*}
那么，我们得到
\begin{equation}
  u_P\sim\int_{-w}^{\infty}\mathrm{e}^{\mathrm{i}\eta^2}\,\mathrm{d}\eta,
\end{equation}
其中
\begin{equation}
  w=d\sqrt{\frac{kD_q}{2D_p(D_q+D_p)}}.
\end{equation}

因此，$P$ 点的强度 $I$ 是
\begin{equation}
  I=\frac{I_0}{2}\left|\sqrt{\frac{2}{\pi}}
  \int_{-w}^{\infty}\mathrm{e}^{\mathrm{i}\eta^2}\,\mathrm{d}\eta\right|^2
  =\frac{I_0}{2}\left[\left(C(w^2)+\frac{1}{2}\right)^2
  +\left(S(w^2)+\frac{1}{2}\right)^2\right],
\end{equation}
其中
\begin{equation*}
  C(z)=\sqrt{\frac{2}{\pi}}\int_0^{\sqrt{z}}\cos\eta^2\,\mathrm{d}\eta,\qquad
  S(z)=\sqrt{\frac{2}{\pi}}\int_0^{\sqrt{z}}\sin\eta^2\,\mathrm{d}\eta
\end{equation*}
称为菲涅耳积分.公式（60.5）解决了所提出的问题，即确定作为 $d$ 的函数的光强度.我们可以看出，$I_0$ 是在照明区内同影子边缘不太接近的点上之强度；更正确一些说，是 $w\gg 1$ 的地方的光强（在 $w\to\infty$ 时 $C(\infty)=S(\infty)=1/2$）.

几何影区与负 $w$ 相应.在 $w$ 为负值且绝对值很大的情况下，容易求出 $I(w)$ 的渐近式.为此，我们按下面方式进行.作分部积分，我们有
\begin{equation*}
  \int_{|w|}^{\infty}\mathrm{e}^{\mathrm{i}\eta^2}\,\mathrm{d}\eta
  =-\frac{1}{2\mathrm{i}|w|}\,\mathrm{e}^{\mathrm{i}w^2}
  +\frac{1}{2\mathrm{i}}\int_{|w|}^{\infty}\mathrm{e}^{\mathrm{i}\eta^2}
  \frac{\mathrm{d}\eta}{\eta^2}.
\end{equation*}
在方程式的右边再作一次分部积分，并重复这种过程，我们得到 $1/|w|$ 的幂级数：
\begin{equation}
  \int_{|w|}^{\infty}\mathrm{e}^{\mathrm{i}\eta^2}\,\mathrm{d}\eta
  =\mathrm{e}^{\mathrm{i}w^2}\left(-\frac{1}{2\mathrm{i}|w|}
  +\frac{1}{4|w|^3}-\cdots\right).
\end{equation}

虽然这样的一个无穷级数不收敛，但是因为当 $|w|$ 很大时，以后的项衰减得很快，因而当 $|w|$ 足够大时第一项已经很好地代表左边函数了（这种级数称为渐近的）.因此，对于强度 $I(w)$，（60.5），我们得到适用于 $w$ 为绝对值很大的负值的情形的渐近公式如下：
\begin{equation}
  I=\frac{I_0}{4\pi w^2}.
\end{equation}

我们看出，在几何影区内，距影边缘很远的地方，强度与到影边缘的距离的平方成反比地趋近于零.

我们现在来考虑 $w$ 的正值，即考虑 $xy$ 平面上方的区域.我们写出
\begin{equation*}
  \int_{-w}^{\infty}\mathrm{e}^{\mathrm{i}\eta^2}\,\mathrm{d}\eta
  =\int_{-\infty}^{+\infty}\mathrm{e}^{\mathrm{i}\eta^2}\,\mathrm{d}\eta
  -\int_{-\infty}^{-w}\mathrm{e}^{\mathrm{i}\eta^2}\,\mathrm{d}\eta
  =(1+\mathrm{i})\sqrt{\frac{\pi}{2}}
  -\int_{w}^{\infty}\mathrm{e}^{\mathrm{i}\eta^2}\,\mathrm{d}\eta.
\end{equation*}
对于足够大的 $w$，可以用方程式右边的积分的渐近式，于是我们有
\begin{equation}
  \int_{-w}^{\infty}\mathrm{e}^{\mathrm{i}\eta^2}\,\mathrm{d}\eta
  \approx(1+\mathrm{i})\sqrt{\frac{\pi}{2}}
  +\frac{1}{2\mathrm{i}w}\,\mathrm{e}^{\mathrm{i}w^2}
\end{equation}
将上式代入（60.5），则我们得到
\begin{equation}
  I=I_0\left[1+\sqrt{\frac{1}{\pi}}\,
  \frac{\sin(w^2-\pi/4)}{w}\right].
\end{equation}

因此，在照明区域内距影边缘甚远处，强度有无穷多的极大值和极小值，因而比值 $I/I_0$ 在1的两边振荡不已.随着 $w$ 增大，振荡的振幅随着与影边缘的距离成反比地减小，极大值与极小值的位置逐渐地彼此接近.

对于小的 $w$，定性来说，$I(w)$ 也有同样的特性.图12就表出了这种特性.在几何影区内，当我们从影的边界移开时，强度单调地减小（在影区边界本身上，$I/I_0=1/4$）.对于正的 $w$，强度将有交替的极大值和极小值.在第一个（最大的）极大值，$I/I_0=1.37$.

\par\nopagebreak\noindent\begin{minipage}{\linewidth}\centering\includegraphics[width=9cm]{fig12.pdf}\\[0.3em]{\small 图 12}\end{minipage}\par\nopagebreak

\section{夫琅禾费衍射}

物理应用中特别感兴趣的衍射现象，是那些在平面平行光束入射到屏上时所发生的现象.由于衍射的结果，光束不再平行，有光线沿着与初始方向不同的方向传播.我们来考虑决定在屏后远距离处衍射光强随方向分布的问题（此处所阐述的问题即为夫琅禾费衍射）.这里我们再次限于同几何光学偏离很小的情形，即假设光线同初始方向的偏离角（衍射角）很小.

这个问题的解决可以从通式（59.2）出发，然后过渡到光源和观察点都离屏无穷远的极限情况.我们正在考虑的情况的特征是，在决定衍射光强度的积分中，取积分的整个波面都是重要的（与此对照，在菲涅尔衍射的情况下，只有屏边缘附近的波面才重要）\footnote{回到公式（60.2）并将其应用到（例如）宽度为 $a$ 的狭缝（而不是孤立屏的边缘），我们容易得到菲涅尔衍射和夫琅禾费衍射的判据.（60.2）中对 $\mathrm{d}z$ 的积分则应当从 0 到 $a$ 进行.菲涅尔衍射对应的情况是，被积函数指数中含 $z^2$ 的项是重要的，且积分上限换为 $\infty$.对于这种情况，我们有
\begin{equation*}
  ka^2\left(\frac{1}{D_p}+\frac{1}{D_q}\right)\gg 1.
\end{equation*}
另一方面，如果把这个不等式反过来，含 $z^2$ 的项可以略去，这种情况就对应于夫琅禾费衍射.}.

然而，较简单的办法是重新处理这个问题，而不借助通式（59.2）.

我们用 $u_0$ 表示假如几何光学严格成立的话屏后可能存在的场.这个场是平面波，但其截面上有某个区域（相应于不透明屏的“阴影”）场为零.用 $S$ 表示该平面截面上场 $u_0$ 不等于零的部分；因为每个这样的平面是平面波的波面，所以在整个 $S$ 面上 $u_0=\mathrm{const}$.

然而，实际上，截面积有限的波不可能是严格的平面波（见 §58）.在它的空间傅里叶展开中，会出现具有不同方向波矢的分量，这正是衍射的起源.

将场 $u_0$ 对波横截面平面内的坐标 $y,z$ 展开为二维傅里叶积分.对于傅里叶分量，我们有：
\begin{equation}
  u_q=\iint u_0\,\mathrm{e}^{-\mathrm{i}\boldsymbol{q}\cdot\boldsymbol{r}}\,
  \mathrm{d}y\,\mathrm{d}z,
\end{equation}
式中矢量 $\boldsymbol{q}$ 是 $yz$ 平面内的常矢量；积分实际上只沿 $yz$ 平面内 $u_0$ 不等于零的 $S$ 部分进行. 如果 $\boldsymbol{k}$ 是入射波的波矢，场分量 $u_q\mathrm{e}^{\mathrm{i}\boldsymbol{q}\cdot\boldsymbol{r}}$ 就给出波矢 $\boldsymbol{k}'=\boldsymbol{k}+\boldsymbol{q}$.因此矢量 $\boldsymbol{q}=\boldsymbol{k}'-\boldsymbol{k}$ 决定了衍射中波矢的变化.因为绝对值 $k=k'=\omega/c$，故 $xy$ 和 $xz$ 平面内的小衍射角 $\theta_y,\theta_z$ 与矢量 $\boldsymbol{q}$ 分量有如下关系：
\begin{equation}
  q_y=\frac{\omega}{c}\theta_y,\qquad
  q_z=\frac{\omega}{c}\theta_z.
\end{equation}

因为同几何光学偏离很小，可以假设场 $u_0$ 展开式中的分量与实际衍射光场的分量恒等，所以公式（61.1）完全解决了我们的问题.

衍射光的强度分布由作为矢量 $\boldsymbol{q}$ 的函数的平方 $|u_q|^2$ 给出.与入射光强度的关系由如下公式决定
\begin{equation}
  \iint u_0^2\,\mathrm{d}y\,\mathrm{d}z
  =\iint|u_q|^2\,\frac{\mathrm{d}q_y\,\mathrm{d}q_z}{(2\pi)^2}
\end{equation}
（比较（49.8））.由此可见，衍射到立体角 $\mathrm{d}o=\mathrm{d}\theta_y\,\mathrm{d}\theta_z$ 内的相对强度由下式给出
\begin{equation}
  \frac{|u_q|^2}{u_0^2}\,\frac{\mathrm{d}q_y\,\mathrm{d}q_z}{(2\pi)^2}
  =\left(\frac{\omega}{2\pi c}\right)^2
  \left|\frac{u_q}{u_0}\right|^2\mathrm{d}o.
\end{equation}

现在我们来考虑两个“互补屏”的夫琅禾费衍射.第一个屏上有一些孔，第二个屏在第一个屏的孔处是不透明的，反之则反.我们用 $u^{(1)}$ 和 $u^{(2)}$ 来代表这些屏所衍射的光的场（两种情况的入射光相同）.因为 $u_q^{(1)}$ 和 $u_q^{(2)}$ 是用在这些屏孔径面上的积分（61.1）来表示的，而两个互补屏上的孔径结合起来就构成整个平面，所以和 $u_q^{(1)}+u_q^{(2)}$ 是屏不存在时场的傅里叶分量，也就是说，它就是入射光.但入射光是具有确定传播方向的严格平面波，所以对于 $\boldsymbol{q}$ 的所有非零值，$u_q^{(1)}+u_q^{(2)}=0$.因此我们有 $u_q^{(1)}=-u_q^{(2)}$，或者对于相应的光强度有
\begin{equation}
  |u_q^{(1)}|^2=|u_q^{(2)}|^2,\qquad
  \text{对于 }\boldsymbol{q}\neq 0.
\end{equation}

这意味着互补的两个屏产生同样的衍射光强度分布（即所谓巴比涅原理）.

这里请注意巴比涅原理的一个有趣的推论.我们来考虑任何一个黑体，即一个吸收所有落于其上光线的物体.按照几何光学，当这个物体被照射时，在它的后面产生一个几何阴影区，其横截面积就等于物体在垂直于入射光线方向的面积.然而，由于衍射的存在，从物体近旁经过的光有一部分偏离了初始方向.结果，在物体后面的远处就不存在完全的阴影了，除了在原方向传播的光以外，还有一部分光在与原方向成一个小角的一些方向行进.求出这种散射光的强度是容易的.为此，我们指出，按照巴比涅原理，由所研究物体的衍射而偏折的光量，等于由屏孔的衍射而偏折的光量，这个孔是从一个不透明的屏切割出来的，孔的形状和面积同物体的横截面一样.但是在孔的夫琅禾费衍射中，所有经过孔的光都偏折了.由此可以断定，被黑体散射的总光量等于落于其表面并被它吸收了的光量.

\begin{xiti}

\noindent{\heiti 习题1}\quad 计算垂直入射在无限长狭缝（宽度为 $2a$）上的平面波的夫琅禾费衍射，狭缝的平行边在不透明的屏上切出.

解：我们选择狭缝的平面为 $yz$ 平面，$z$ 轴沿着狭缝（图13表示该屏的截面）.对于垂直入射光，狭缝的平面是波面之一. 我们选择它作为（61.1）中积分的面. 因为狭缝是无限长的，光只在 $xy$ 平面内偏折（积分（61.1）对于 $q_z\neq 0$ 等于零）.

\par\nopagebreak\noindent\begin{minipage}{\linewidth}\centering\includegraphics[width=4.5cm]{fig13.pdf}\\[0.3em]{\small 图 13}\end{minipage}\par\nopagebreak

因而场只应当在 $y$ 坐标中展开：
\begin{equation*}
  u_q=u_0\int_{-a}^{a}\mathrm{e}^{-\mathrm{i}qy}\,\mathrm{d}y
  =\frac{2u_0}{q}\sin qa.
\end{equation*}

在角范围 $\mathrm{d}\theta$ 内衍射光的强度是
\begin{equation*}
  \mathrm{d}I=\frac{I_0}{2a}\left|\frac{u_q}{u_0}\right|^2
  \frac{\mathrm{d}q}{2\pi}
  =\frac{I_0}{\pi ak}\,\frac{\sin^2 ka\theta}{\theta^2}\,\mathrm{d}\theta,
\end{equation*}
式中，$k=\omega/c$，$I_0$ 是入射到狭缝上的光的总强度.

$\mathrm{d}I/\mathrm{d}\theta$ 作为衍射角的函数具有图14所显示的形状.随着 $\theta$ 从 $\theta=0$ 朝两边增加，强度历经一系列高度迅速减小的极大值. 相邻的极大值被处于 $\theta=n\pi/ka$ 点的极小值隔开（式中 $n$ 是整数）；在极小值处，强度下降到零.

\par\nopagebreak\noindent\begin{minipage}{\linewidth}\centering\includegraphics[width=7cm]{fig14.pdf}\\[0.3em]{\small 图 14}\end{minipage}\par\nopagebreak

\noindent{\heiti 习题2}\quad 计算衍射光栅的夫琅禾费衍射. 衍射光栅是一个平面屏上面刻了一系列完全一样的平行狭缝，狭缝的宽度为 $2a$，相邻狭缝之间不透明屏的宽度为 $2b$，狭缝的总数为 $N$.

解：我们选定光栅的平面作为 $yz$ 平面，$z$ 轴与狭缝平行.衍射只发生在 $xy$ 平面内，将（61.1）积分得到：
\begin{equation*}
  u_q=u_q'\sum_{n=0}^{N-1}\mathrm{e}^{-2\mathrm{i}nqd}
  =u_q'\,\frac{1-\mathrm{e}^{-2\mathrm{i}Nqd}}{1-\mathrm{e}^{-2\mathrm{i}qd}},
\end{equation*}
式中，$d=a+b$，而 $u_q'$ 则是沿一条狭缝积分的结果.利用习题1的结果，我们得到：
\begin{equation*}
  \mathrm{d}I=\frac{I_0a}{N\pi}\left(\frac{\sin Nqd}{\sin qd}\right)^2
  \left(\frac{\sin qa}{qa}\right)^2\mathrm{d}q
  =\frac{I_0}{N\pi ak}\left(\frac{\sin Nk\theta d}{\sin k\theta d}\right)^2
  \frac{\sin^2 ka\theta}{\theta^2}\,\mathrm{d}\theta
\end{equation*}
（$I_0$ 是经过所有狭缝的总光强度）.

在有很多狭缝的情形下，就是当 $N\to\infty$ 时，这个公式可以写成另一形式. 对于 $q=\pi n/d$ 的值（$n$ 是整数），$\mathrm{d}I/\mathrm{d}q$ 有一极大值；在这些极大值的近旁（即 $qd=n\pi+\varepsilon$，$\varepsilon$ 很小）：
\begin{equation*}
  \mathrm{d}I=I_0a\left(\frac{\sin qa}{qa}\right)^2
  \frac{\sin^2 N\varepsilon}{\pi N\varepsilon^2}\,\mathrm{d}q.
\end{equation*}

但是当 $N\to\infty$ 时，下面的公式成立\footnote{对于 $x\neq 0$，方程左边的函数为零，而按照傅里叶级数理论中熟知的公式
\begin{equation*}
  \lim_{N\to\infty}\left(\frac{1}{\pi}\int_{-a}^{a}f(x)\frac{\sin^2 Nx}{Nx^2}\,\mathrm{d}x\right)=f(0).
\end{equation*}
由此可见，函数的特性实际上与 $\delta$ 函数的特性相合（见 §28 的脚注）.}：
\begin{equation*}
  \lim_{N\to\infty}\frac{\sin^2 Nx}{\pi Nx^2}=\delta(x).
\end{equation*}

% ---- 脚注①（原书 p.176）已由脚注代理对照原书扫描页转录排入（2026-10-10）；
% ---- 原书中“由此可见……相合（见 §28 的脚注）”一句属于本脚注，已从正文移入。 ----

因而在每个极大值近旁，我们得到：
\begin{equation*}
  \mathrm{d}I=I_0\,\frac{a}{d}\left(\frac{\sin qa}{qa}\right)^2
  \delta(\varepsilon)\,\mathrm{d}\varepsilon,
\end{equation*}
就是说，在极大值宽度无限窄的极限下，第 $n$ 个极大值的总光强是
\begin{equation*}
  I^{(n)}=I_0\,\frac{d}{\pi^2 a}\,
  \frac{\sin^2(n\pi a/d)}{n^2}.
\end{equation*}

\noindent{\heiti 习题3}\quad 光垂直射入半径为 $a$ 的圆孔径的平面，求衍射光强度随方向的分布.

解：我们引入柱坐标 $z,r,\varphi$，$z$ 轴经过圆孔径的中心并垂直于圆孔径的平面. 显然衍射是绕 $z$ 轴对称的，所以矢量 $\boldsymbol{q}$ 只有径向分量 $q_r=q=k\theta$.从方向 $\boldsymbol{q}$ 测量角 $\varphi$，在（61.1）中沿孔径平面进行积分，我们得到：
\begin{equation*}
  u_q=u_0\int_0^{a}\int_0^{2\pi}
  \mathrm{e}^{-\mathrm{i}qr\cos\varphi}\,r\,\mathrm{d}\varphi\,\mathrm{d}r
  =2\pi u_0\int_0^{a}\mathrm{J}_0(qr)\,r\,\mathrm{d}r,
\end{equation*}
式中，$\mathrm{J}_0$ 是零阶的贝塞耳函数，用熟知的公式
\begin{equation*}
  \int_0^{a}\mathrm{J}_0(qr)\,r\,\mathrm{d}r
  =\frac{a}{q}\,\mathrm{J}_1(aq),
\end{equation*}
我们于是得到
\begin{equation*}
  u_q=\frac{2\pi u_0a}{q}\,\mathrm{J}_1(aq),
\end{equation*}
按照（61.4），我们得出衍射到立体角 $\mathrm{d}o$ 内光的强度
\begin{equation*}
  \mathrm{d}I=I_0\,\frac{\mathrm{J}_1^2(ak\theta)}{\pi\theta^2}\,\mathrm{d}o,
\end{equation*}
式中，$I_0$ 是入射到孔径上光的总强度.

\end{xiti}
